% Generated draft: requires mathematical and source review.
\documentclass[11pt]{article}
\usepackage[margin=25mm]{geometry}
\usepackage{fontspec}
\setmainfont{Latin Modern Roman}
\usepackage{amsmath,amssymb,mathtools,unicode-math}
\setmathfont{Latin Modern Math}
\usepackage{enumitem,needspace,xurl,hyperref}
\setlength{\parindent}{0pt}
\setlength{\parskip}{5pt}
\setlength{\emergencystretch}{3em}
\newcommand{\N}{\mathbb{N}}
\newcommand{\Z}{\mathbb{Z}}
\newcommand{\Q}{\mathbb{Q}}
\newcommand{\R}{\mathbb{R}}
\newcommand{\C}{\mathbb{C}}
\newcommand{\F}{\mathbb{F}}
\newcommand{\A}{\mathbb{A}}
\newcommand{\PP}{\mathbb{P}}
\newcommand{\E}{\mathbb{E}}
\newcommand{\eps}{\varepsilon}
\newcommand{\dd}{\,\mathrm{d}}
\begin{document}
\section*{UnsolvedMath ID: 801164}
\textbf{Source problem code:} catalogue entry
\subsection{Definitions and mathematical statement}
The supplied excerpt does not define its wavelet system or its realization of \(H^2(\R)\). Under the standard dyadic, integer-translation convention and the positive-frequency realization of the Hardy space, the question has the following precise formulation. These conventions require confirmation against the cited source.

Use the Fourier transform \(\widehat f(\xi)=\int_{\R}f(x)e^{-2\pi i x\xi}\,\dd x\), extended to \(L^2(\R;\C)\), and set
\[
H^2(\R)=\{f\in L^2(\R;\C):\widehat f(\xi)=0\text{ for almost every }\xi<0\}.
\]
The Schwartz class is
\[
\mathcal S(\R)=\left\{f\in C^\infty(\R;\C):\sup_{x\in\R}|x^m f^{(n)}(x)|<\infty\text{ for all }m,n\in\Z_{\geq0}\right\}.
\]
For \(\psi\in H^2(\R)\), define
\[
\psi_{j,k}(x)=2^{j/2}\psi(2^j x-k),\qquad j,k\in\Z.
\]
This family is a Riesz basis of \(H^2(\R)\) if its closed linear span is \(H^2(\R)\) and there exist constants \(0<A\leq B<\infty\) such that, for every finitely supported family \((c_{j,k})_{(j,k)\in\Z^2}\) of complex numbers,
\[
A\sum_{j,k\in\Z}|c_{j,k}|^2
\leq
\left\|\sum_{j,k\in\Z}c_{j,k}\psi_{j,k}\right\|_{L^2(\R)}^2
\leq
B\sum_{j,k\in\Z}|c_{j,k}|^2.
\]
Equivalently, a Riesz basis is the image of an orthonormal basis under a bounded invertible linear operator on the Hilbert space.

\textbf{Question, with the conventions above.} Does there exist
\[
\psi\in\mathcal S(\R)\cap H^2(\R)
\]
such that \((\psi_{j,k})_{j,k\in\Z}\) is a Riesz basis of \(H^2(\R)\)?

The source also defines Riesz wavelets for \(L^2(\R)\), but the question here specifically concerns \(H^2(\R)\). Every Riesz basis has a dual Riesz basis; the source emphasizes that the dual of a wavelet Riesz basis need not itself be a wavelet system. When it is, the source calls the wavelet biorthogonal. The question above imposes no such additional requirement on the dual.
\subsection{Short English statement}
Does the real-line Hardy space admit a wavelet Riesz basis generated by a Schwartz function? The displayed formulation uses the standard dyadic wavelet system and positive-frequency Hardy space, conventions that the supplied excerpt leaves unspecified. No requirement that the dual basis also be a wavelet system is imposed.
\subsection{Sources}
\begin{itemize}
\item Local UnsolvedMath registry, permanent ID \(801164\), exact source problem code \texttt{catalogue entry}, titled \emph{Open Problems in Wavelet Theory --- Problem 3}. The supplied registry labels this problem \emph{open}; that label is supplied curation, not an independently verified current literature assessment.
\item Bownik and Rzeszotnik, \emph{Open Problems in Wavelet Theory} (2019), Problem 3, \url{https://arxiv.org/abs/1905.08230}. The supplied provenance records extraction from arXiv LaTeX and an access date of 2026-10-16. This access information is supplied provenance; no independent source visit or verification is asserted.
\item The excerpt attributes the existence problem to Seip through the reference label [Seip]. No bibliographic entry for [Seip] is supplied, so the original reference and its precise scope remain unverified. The registry separately lists Bownik and Rzeszotnik as proposers and 2019 as the proposed year.
\end{itemize}
\end{document}
